\documentclass[10pt,a4paper]{amsart}
\usepackage[margin=19mm]{geometry}
\usepackage{amsmath,amssymb,mathtools}
\usepackage[scr=rsfs,cal=euler]{mathalfa}
\usepackage{xfrac}
\usepackage[unicode,hidelinks,bookmarks=false]{hyperref}
\newcommand{\ie}{i.e.\ }
\newcommand{\cf}{cf.\ }
\newcommand{\resp}{respectively\ }
\newcommand{\Subsection}{\subsection}
\providecommand{\nobreakdash}{\nobreak\hbox{-}}
\setlength{\emergencystretch}{2em}
\multlinegap0pt
\usepackage{bm}
\usepackage{bbm}
\usepackage{dsfont}
\usepackage{xspace}
\usepackage{cancel}
\usepackage{mathtools}
\usepackage{amsthm}
\makeatletter
\@ifundefined{lemma}{\newtheorem{lemma}{Lemma}}{}
\@ifundefined{prop}{\newtheorem{prop}[lemma]{Proposition}}{}
\makeatother
\newcommand{\cG}{\mathcal{G}}
\renewcommand{\phi}{\varphi}
\title[Bounded cohomological induction for transverse measured grou]{Bounded cohomological induction for transverse measured groupoids}
\author{Tobias Hartnick, Filippo Sarti}
\date{Source: arXiv:2510.20656v1 (CC BY 4.0)}
\begin{document}
\maketitle

\noindent\emph{Source and licence.} Bounded cohomological induction for transverse measured groupoids. Version arXiv:2510.20656v1 (CC BY 4.0), distributed under the Creative Commons Attribution 4.0 International licence, as recorded in the arXiv licence metadata for this exact version and shown on the abstract page https://arxiv.org/abs/2510.20656v1. Full rights notice: licenses/arxiv-2510.20656-CC-BY-4.0.txt.\par\smallskip

\noindent\textit{Editorial scope.} Counted unit: the Prop whose source label is \texttt{psiIso} in the article (source0.tex lines 1199--1205, source label \texttt{psiIso}), delivered together with the article's own complete original proof of it (source0.tex lines 1206--1246, one \texttt{proof} environment of 41 lines). No mathematics is written, repaired or re-derived: statement and proof are the article's own text line for line. Required context retained uncounted: FiberProdRed (lines 1141--1143); AstAction1 (lines 1159--1161). Every definition, display and label of those lines is retained unchanged. Omitted from this package and not redistributed: 1--1140, 1144--1158, 1162--1198, 1247--1797. Nothing inside a retained excerpt is omitted or altered: every retained body is delivered exactly as the article's lines are written, so no omission receipt of any kind accompanies this package. Citations: the retained excerpts carry no citation command at all, so no bibliography entry of the article is transcribed and no reference is redistributed. Reference adaptation: none was needed and none was made; every reference inside the delivered unit resolves inside it, so no unresolved reference or citation is produced. Printed slips kept literal: no printed word, symbol, hypothesis or number of the counted statement or of its proof was altered. Layout and class: the article's own class is \texttt{amsart} with packages including babel, enumerate, graphics, color, pgf, epsfig, inputenc, xy, tikz-cd, bbm, comment, amssymb, amsmath, amsthm; the wrapper is the lane's pinned \texttt{amsart} editorial wrapper at 10pt on A4, which restates the theorem-like environments the retained text needs and adds the article's own macro definitions that the retained text needs (\texttt{\textbackslash cG}, \texttt{\textbackslash phi}), with extra wrapper packages (bm, bbm, dsfont, xspace, cancel, mathtools). The delivered numbering is the unit's own. Assets: no retained span contains a figure environment, a graphics-inclusion command, a PDF-page inclusion, a table or a TikZ picture; no image file is copied into this package, the package declares no asset dependency and no image file is redistributed with it. Licence: version arXiv:2510.20656v1 of this article is distributed under the Creative Commons Attribution 4.0 International licence, as recorded in the arXiv licence metadata for this exact version and shown on the abstract page https://arxiv.org/abs/2510.20656v1. A full rights notice is delivered at licenses/arxiv-2510.20656-CC-BY-4.0.txt. Characters: every retained excerpt is pure ASCII, so the delivered engine, XeLaTeX with its default Latin Modern text font, reports no missing character.
\begin{equation}\label{FiberProdRed}
\cG^{(1)}\times_{Y} X = \{[h,x] \mid h \in G, x \in X,  h\beta(x)x \in Y\}.
\end{equation}

\begin{equation}\label{AstAction1}
((g, y), k) \ast [h,x] \coloneqq ([gh\sigma(k,x)^{-1}, kx]). \tag{$\ast$}
\end{equation}

\begin{prop}\label{psiIso} There is a well-defined Borel isomorphism
\begin{equation}\label{psi}
  \psi: G\times Y\to \mathcal{G}^{(1)} \times_Y X\,,\;\;\; \psi(\gamma,y)=[\gamma^{-1}\beta(\gamma y)^{-1}, \gamma y]\end{equation}
 with inverse given by
  $$\phi: \mathcal{G}^{(1)} \times_Y X \to G\times Y\,,\;\;\; \phi([h,x])=(\beta(x)^{-1} h^{-1} , h \beta(x) x).$$
 Moreover, this isomorphism intertwines the induced $(\cG \times G)$-action action on $G\times Y$ with the action on the transfer space given by \eqref{AstAction1}.
\end{prop}

\begin{proof} We first show that $\psi$ is well-defined. Since $s(\beta(\gamma y), \gamma y) = t(\gamma, y)$ we can form the product
\[
(\beta(\gamma y), \gamma y)(\gamma, y) = (\beta(\gamma y)\gamma, y),
\]
which has source $y$ and target $\beta(\gamma y)\gamma y \in Y$, hence is contained in $\cG^{(1)}$. Now pass to the inverse to see that $((\beta(\gamma y), \gamma, y)(\gamma, y))^{-1} \in \cG^{(1)}$; this shows that $\psi$ is well-defined, and it is Borel as a composition of Borel maps. Similarly, $\phi$ is well-defined by \eqref{FiberProdRed} and Borel by definition. We compute 
  \begin{align*}
    \psi \circ \phi([h,x]) &= \psi(\beta(x)^{-1} h^{-1},h \beta(x) x)\\
     & = [h \beta(x) \beta(h \beta(x)\beta(x)^{-1} h^{-1}x)^{-1} , \beta(x)^{-1} h^{-1}h\beta(x) x ] = [h,x]
  \end{align*}
  and dually
  \begin{align*}
  \phi\circ \psi(\gamma,y)&= \phi ([\gamma^{-1}\beta( \gamma y)^{-1},\gamma y])\\
 &= (\beta(\gamma y)^{-1}\beta( \gamma y) \gamma , \gamma^{-1}\beta(\gamma y)^{-1} \beta(\gamma y)\gamma y)=(\gamma,y),
  \end{align*}
  which shows that $\psi$ and $\phi$ are mutually inverse Borel isomorphisms. To see equivariance we observe that, on the one hand,
\begin{align*}
    \psi(((g,y),k)\star (\gamma,y))&=\psi(k\gamma g^{-1},gy) =([g\gamma^{-1}k^{-1} \beta(k\gamma y)^{-1},k\gamma y]),
\end{align*}
and on the other hand,
\begin{align*}
   ((g,y),k)  \ast \psi(\gamma,y)&=   ((g,y),k)\ast  ([\gamma^{-1}\beta(\gamma y)^{-1},\gamma y])\\
   &=[g(\gamma^{-1}\beta(\gamma y)^{-1})\sigma(k,\gamma y)^{-1},k\gamma y].
\end{align*}
Spelling out the definition of $\sigma$ now yields
\[
g(\gamma^{-1}\beta(\gamma y)^{-1})\sigma(k,\gamma y)^{-1} = g\gamma^{-1}\beta(\gamma y)^{-1}(\beta(k \gamma y)k \beta(\gamma y)^{-1})^{-1}
= g \gamma^{-1}k^{-1}\beta(k\gamma y)^{-1},
\]
which shows that the two expressions are equal. This concludes the proof.
% Similarly, we have 
% \begin{align*}
%   \varphi(((g,y),k)\ast[h,x])&=\varphi(gh\sigma(k,x)^{-1},kx)\\
%   &=(\beta(kx)^{-1}\sigma(k,x)h^{-1}g^{-1},gh\sigma(k,x)^{-1}\beta(kx)kx)
% \end{align*}
% and 
% \begin{align*}
%   ((g,y),k)\star \varphi([h,x])&=((g,y),k)\star(\beta(x)^{-1}h^{-1},h\beta(x)x)\\
%   &=(k\beta(x)^{-1}h^{-1}g^{-1},gh\beta(x)x)\,,
% \end{align*}
% so we can conclude using again the definition of $\sigma$. 
\end{proof}

\end{document}
